% Executive summary of how a register spill is defined and measured.
% Build:  tectonic --outdir benchmarks benchmarks/spill_definition.tex
% (or:    pdflatex -output-directory benchmarks benchmarks/spill_definition.tex)
\documentclass[10pt,a4paper]{article}
\usepackage[a4paper,top=16mm,bottom=16mm,left=20mm,right=20mm]{geometry}
\usepackage{booktabs}
\usepackage{tabularx}
\usepackage{lmodern}
\usepackage[T1]{fontenc}
\usepackage{amsmath}
\usepackage{microtype}
\usepackage{titlesec}

\pagestyle{empty}
\setlength{\parindent}{0pt}
\setlength{\parskip}{5pt}
\linespread{1.06}\selectfont
\raggedbottom

\titleformat{\section}{\normalfont\large\bfseries}{\thesection.}{0.6em}{}
\titlespacing*{\section}{0pt}{12pt plus 2pt minus 2pt}{5pt}

% Author-block markup kept in the IEEEtran spelling the paper uses, so the same
% block can be moved between documents unchanged. This report stays
% single-column: the rule tables need the full text width to stay legible.
% The affiliation block contains its own \\ line breaks. article's \author
% puts each author in a tabular, where a brace group cannot span a row break,
% so the affiliation goes into a nested tabular instead.
\newcommand{\IEEEauthorblockN}[1]{{\normalsize #1}}
\newcommand{\IEEEauthorblockA}[1]{%
  \\[2pt]{\small\begin{tabular}[t]{@{}c@{}}#1\end{tabular}}}

\newcommand{\LongDate}{%
    \number\day\space
    \ifcase\month
    \or January%
    \or February%
    \or March%
    \or April%
    \or May%
    \or June%
    \or July%
    \or August%
    \or September%
    \or October%
    \or November%
    \or December%
    \fi
    \space \number\year
}

% X columns whose share of the width can be tuned, so "Test" gets the room it
% needs and "Rationale" is never squeezed.
\newcolumntype{L}[1]{>{\raggedright\arraybackslash\hsize=#1\hsize}X}
\newcolumntype{R}{>{\ttfamily}l}

\setlength{\tabcolsep}{6pt}
\renewcommand{\arraystretch}{1.28}

\newcommand{\tablehead}[2]{%
  \nobreak{\small\textbf{Table #1.} #2}\par\nobreak\vspace{2pt}}

\title{
    \Large Improving CPU Efficiency for Data-Intensive Applications \\
    \Large through Cache-Aware Spill Optimization \\[6pt]
    \large Executive Summary: How a Spill Is Defined and Measured \\[5pt]
    \normalsize gem5 TimingSimpleCPU $\cdot$ RISC-V (RV64GC)
    $\cdot$ SPEC CPU2017 $\cdot$ metric schema 7
}

\author{
    \IEEEauthorblockN{L\"utfullah Burak Kaya}
    \IEEEauthorblockA{
        Ozyegin University \\
        burak.kaya.50711@ozu.edu.tr
    }
    \and
    \IEEEauthorblockN{Ismail Akturk}
    \IEEEauthorblockA{
        Ozyegin University \\
        ismail.akturk@ozyegin.edu.tr
    }
}

\date{\small \LongDate}

\begin{document}

\maketitle
\thispagestyle{empty}
\vspace{-9mm}

\section{What is measured}

A spill is detected \emph{dynamically}, from the memory behaviour of the
running program, not from compiler metadata. The detector watches every
committed load and store and looks for the signature of compiler spill code: a
value written to a stack slot through the stack or frame pointer, and later
read back from that slot.

All counting is confined to the region of interest (ROI), marked by
\texttt{m5\_work\_begin} and \texttt{m5\_work\_end} in the benchmark sources.
State is cleared at ROI entry, so no pair straddles the boundary.

\section{Store side}

Three tests decide whether a committed store is recorded as a spill write.

\tablehead{1}{Store-side rules, applied when the store commits.}

\begin{tabularx}{\textwidth}{@{}l L{1.32} L{0.68}@{}}
\toprule
& \textbf{Test} & \textbf{Rationale} \\
\midrule
S1 & The store executes inside the ROI
   & Measurement window \\
S2 & The target address lies in the stack region,
     $\texttt{stackMin} \le \text{addr} < \texttt{stackBase}$
   & Spill slots live on the stack \\
S3 & The address base register is \texttt{sp}~(\texttt{x2})
     or \texttt{s0}/\texttt{fp}~(\texttt{x8})
   & How spill code addresses the stack \\
\bottomrule
\end{tabularx}

A passing store is counted in \texttt{spill\_store}, tagged for the cache
statistics, and remembered in a map keyed by its address. A later store to that
address replaces the entry, beginning a new \emph{generation}.

\section{Load side}

Six tests decide whether a committed load reads back a recorded spill.

\tablehead{2}{Load-side rules, applied when the load commits.}

\begin{tabularx}{\textwidth}{@{}l L{1.32} L{0.68}@{}}
\toprule
& \textbf{Test} & \textbf{Rationale} \\
\midrule
L1 & The load executes inside the ROI
   & Measurement window \\
L2 & The source address lies in the stack region
   & As S2 \\
L3 & A recorded store exists at exactly this address
   & There is a value to read back \\
L4 & The load is later in time than that store
   & A spill is written before it is read \\
L5 & Load and store use the same access width
   & The same value, not an overlapping one \\
L6 & Load and matched store are both base-register relative, per S3
   & Both halves must look like spill code \\
\bottomrule
\end{tabularx}

A load passing all six is counted in \texttt{spill\_load}; the first one for a
generation also completes a \texttt{spill\_event}. Pairs that satisfy L1--L5
but fail L6 are ordinary stack data reuse, counted separately.

\section{Counters}

\tablehead{3}{Reported counters and how each is derived.}

\begin{tabularx}{\textwidth}{@{}R X@{}}
\toprule
\normalfont\textbf{Counter} & \textbf{Definition} \\
\midrule
spill\_store
  & Write side: every store passing S1--S3, read back later or not. \\
spill\_load
  & Read-back side: every load passing L1--L6, repeats included. \\
spill\_event
  & One (\texttt{spill\_store}, first \texttt{spill\_load}) pair: one
    \emph{potential spill}. \\
total\_spill\_insts
  & \texttt{spill\_store} $+$ \texttt{spill\_load}. \\
spill\_store\_unpaired
  & \texttt{spill\_store} $-$ \texttt{spill\_event}: written, never read back
    inside the ROI. \\
repeat\_spill\_load
  & \texttt{spill\_load} $-$ \texttt{spill\_event}: read-backs after the
    first. \\
\bottomrule
\end{tabularx}

The two sides are decided at different moments, so two inequalities always
hold: $\texttt{spill\_event} \le \texttt{spill\_store}$, since a value may be
written and never read back; and
$\texttt{spill\_load} \ge \texttt{spill\_event}$, since a value may be read
back many times.

\section{What this does and does not prove}

The rules establish that a stack slot was written through the stack or frame
pointer and, for a \texttt{spill\_event}, read back at the same width.

They do \emph{not} establish that the register allocator ran out of registers.
Three other populations pass exactly the same tests and cannot be told apart
at the microarchitectural level:

\begin{itemize}
\setlength{\itemsep}{1pt}\setlength{\topsep}{2pt}
\item ABI-mandated callee-save save and restore in prologues and epilogues;
\item ordinary scalar locals the compiler kept in memory;
\item incoming arguments written to their stack home.
\end{itemize}

The figures are therefore an \textbf{upper bound} on true register-allocator
spilling, and should be read as \emph{potential} spills. The write side is the
looser of the two: at store time a later read-back is unknowable, so
\texttt{spill\_store} need only pass S1--S3.

For an exact classification the detector also accepts a compiler-generated PC
map (\texttt{spill\_pc\_map}), replacing tests S3 and L6 with the addresses
the compiler emitted as spill and reload code.

\section{Implementation}

\begin{tabularx}{\textwidth}{@{}X l@{}}
\toprule
\textbf{Component} & \textbf{Source} \\
\midrule
Rules and counters
  & \texttt{src/cpu/simple/spill\_detector.\{hh,cc\}} \\
Base-register test, commit hook
  & \texttt{src/cpu/simple/timing.cc} \\
ROI markers, statistics boundaries
  & \texttt{src/sim/pseudo\_inst.cc} \\
Per-population cache counters
  & \texttt{src/mem/cache/base.\{hh,cc\}} \\
\bottomrule
\end{tabularx}

The rule set is covered by 21 hand-verified RISC-V assembly test cases in
\texttt{tests/spill\_test/}, which also assert that the gem5 cache statistics
and the detector counters agree at both cache levels.

\end{document}
